\documentclass[12pt, a4paper]{article}

% ── Packages ──────────────────────────────────────────────────────────────────
\usepackage[margin=1in]{geometry}
\usepackage{amsmath, amssymb, amsthm}
\usepackage{mathtools}
\usepackage{hyperref}
\usepackage{xcolor}
\usepackage{listings}
\usepackage{enumitem}
\usepackage{booktabs}
\usepackage{titlesec}
\usepackage{fancyhdr}
\usepackage{datetime2}
\usepackage{tcolorbox}
\tcbuselibrary{skins, breakable}
\usepackage[utf8]{inputenc}
\usepackage{newunicodechar}
\newunicodechar{ℝ}{\ensuremath{\mathbb{R}}}
\newunicodechar{ℕ}{\ensuremath{\mathbb{N}}}
\newunicodechar{ℤ}{\ensuremath{\mathbb{Z}}}
\newunicodechar{α}{\ensuremath{\alpha}}
\newunicodechar{β}{\ensuremath{\beta}}
\newunicodechar{γ}{\ensuremath{\gamma}}
\newunicodechar{δ}{\ensuremath{\delta}}
\newunicodechar{ε}{\ensuremath{\varepsilon}}
\newunicodechar{φ}{\ensuremath{\phi}}
\newunicodechar{ω}{\ensuremath{\omega}}
\newunicodechar{Ω}{\ensuremath{\Omega}}
\newunicodechar{∀}{\ensuremath{\forall}}
\newunicodechar{∃}{\ensuremath{\exists}}
\newunicodechar{→}{\ensuremath{\to}}
\newunicodechar{←}{\ensuremath{\leftarrow}}
\newunicodechar{↔}{\ensuremath{\leftrightarrow}}
\newunicodechar{∧}{\ensuremath{\wedge}}
\newunicodechar{∨}{\ensuremath{\vee}}
\newunicodechar{¬}{\ensuremath{\neg}}
\newunicodechar{≠}{\ensuremath{\neq}}
\newunicodechar{≤}{\ensuremath{\leq}}
\newunicodechar{≥}{\ensuremath{\geq}}
\newunicodechar{∈}{\ensuremath{\in}}
\newunicodechar{∉}{\ensuremath{\notin}}
\newunicodechar{⊆}{\ensuremath{\subseteq}}
\newunicodechar{∣}{\ensuremath{\mid}}
\newunicodechar{⟨}{\ensuremath{\langle}}
\newunicodechar{⟩}{\ensuremath{\rangle}}
\newunicodechar{∑}{\ensuremath{\sum}}
\newunicodechar{∫}{\ensuremath{\int}}
\newunicodechar{∞}{\ensuremath{\infty}}
\newunicodechar{∂}{\ensuremath{\partial}}
\newunicodechar{∇}{\ensuremath{\nabla}}
\newunicodechar{Δ}{\ensuremath{\Delta}}
\newunicodechar{⁻}{\ensuremath{^{-}}}
\newunicodechar{⟫}{\ensuremath{\rangle\!\rangle}}
\newunicodechar{⟪}{\ensuremath{\langle\!\langle}}
\newunicodechar{‖}{\ensuremath{\|}}
\newunicodechar{·}{\ensuremath{\cdot}}
\newunicodechar{×}{\ensuremath{\times}}

% ── Colors ────────────────────────────────────────────────────────────────────
\definecolor{leanblue}{RGB}{41, 103, 178}
\definecolor{mathgreen}{RGB}{0, 128, 80}
\definecolor{sorryred}{RGB}{180, 40, 40}
\definecolor{codebg}{RGB}{245, 247, 250}
\definecolor{titlegray}{RGB}{60, 60, 80}

% ── Hyperref setup ────────────────────────────────────────────────────────────
\hypersetup{
  colorlinks=true,
  linkcolor=leanblue,
  urlcolor=leanblue,
  citecolor=mathgreen
}

% ── Lean listings style ───────────────────────────────────────────────────────
\lstdefinelanguage{Lean4}{
  keywords={def, theorem, lemma, structure, class, instance, import,
            noncomputable, variable, section, end, namespace, open,
            where, by, fun, let, have, show, exact, apply, intro,
            rw, simp, calc, sorry, rfl, constructor, infer_instance},
  keywordstyle=\color{leanblue}\bfseries,
  comment=[l]{--},
  commentstyle=\color{gray}\itshape,
  string=[b]",
  stringstyle=\color{mathgreen},
  sensitive=true,
  literate=
    {∀}{{$\forall$}}1
    {∃}{{$\exists$}}1
    {→}{{$\to$}}1
    {←}{{$\leftarrow$}}1
    {↔}{{$\leftrightarrow$}}1
    {∧}{{$\wedge$}}1
    {∨}{{$\vee$}}1
    {¬}{{$\neg$}}1
    {≠}{{$\neq$}}1
    {≤}{{$\leq$}}1
    {≥}{{$\geq$}}1
    {∈}{{$\in$}}1
    {∉}{{$\notin$}}1
    {⊆}{{$\subseteq$}}1
    {ℝ}{{$\mathbb{R}$}}1
    {ℕ}{{$\mathbb{N}$}}1
    {ℤ}{{$\mathbb{Z}$}}1
    {α}{{$\alpha$}}1
    {β}{{$\beta$}}1
    {ε}{{$\varepsilon$}}1
    {δ}{{$\delta$}}1
    {∣}{{$\mid$}}1
    {•}{{$\bullet$}}1
    {:=}{{$:=$}}2
}

\lstset{
  language=Lean4,
  basicstyle=\ttfamily\small,
  backgroundcolor=\color{codebg},
  frame=single,
  framerule=0.4pt,
  rulecolor=\color{leanblue!40},
  xleftmargin=10pt,
  xrightmargin=10pt,
  breaklines=true,
  showstringspaces=false,
  tabsize=2,
  captionpos=b
}

% ── Theorem environments ──────────────────────────────────────────────────────
\theoremstyle{definition}
\newtheorem{definition}{Definition}[section]
\newtheorem{remark}{Remark}[section]
\newtheorem{todo}{TODO}[section]

\theoremstyle{plain}
\newtheorem{theorem}{Theorem}[section]
\newtheorem{lemma}[theorem]{Lemma}

% ── Custom boxes ──────────────────────────────────────────────────────────────
\newtcolorbox{leanbox}[1][]{
  colback=codebg, colframe=leanblue!60,
  fonttitle=\bfseries\small, title=Lean 4,
  #1
}

\newtcolorbox{statusbox}[2][]{
  colback=#2!8, colframe=#2!60,
  fonttitle=\bfseries\small, title=Status,
  #1
}

% ── Header / Footer ───────────────────────────────────────────────────────────
\pagestyle{fancy}
\fancyhf{}
\rhead{\small PDE Formalization in Lean 4}
\lhead{}
\cfoot{\thepage}
\renewcommand{\headrulewidth}{0.4pt}

% ── Title formatting ──────────────────────────────────────────────────────────
\titleformat{\section}
  {\large\bfseries\color{titlegray}}{\thesection}{1em}{}[\titlerule]
\titleformat{\subsection}
  {\normalsize\bfseries\color{leanblue}}{\thesubsection}{1em}{}

% ── Document ──────────────────────────────────────────────────────────────────
\begin{document}

% Title page
\begin{titlepage}
  \centering
  \vspace*{2cm}
  {\Huge\bfseries\color{titlegray} Formalizing PDE Theory\\[0.4em]
   \large in Lean 4 and Mathlib}\\[3cm]
  \begin{tcolorbox}[colback=leanblue!5, colframe=leanblue!40, width=0.85\textwidth]
    \centering
    \textbf{Main Objective}\\[0.5em]
    Introduce and organize foundational PDE objects inside Mathlib
    so that future formalization of major PDE theorems becomes possible.
    The immediate goal is not to prove large theorems, but to build
    the infrastructure.
  \end{tcolorbox}
  \vfill
  {\small Based on: L.C. Evans, \textit{Partial Differential Equations}, AMS 1998}\\
  {\small Reference: \href{https://leanprover-community.github.io/mathematics_in_lean/}
    {Mathematics in Lean}, \href{https://leanprover-community.github.io/mathlib4_docs/}{Mathlib4 Docs}}
\end{titlepage}

\tableofcontents
\newpage

% ══════════════════════════════════════════════════════════════════════════════
\section{Project Overview}
% ══════════════════════════════════════════════════════════════════════════════

\subsection{Philosophy}

The project follows the principle:
\begin{center}
  \textit{Propositions are types. Proofs are terms.}
\end{center}
Every mathematical statement in Lean is a \texttt{Prop} (a type), and proving it
means constructing a term of that type. The compiler \emph{is} the verifier ---
if Lean accepts the code with no errors and no \texttt{sorry} warnings, the proof
is formally and absolutely correct.

\subsection{Golden Rule}

\begin{tcolorbox}[colback=mathgreen!6, colframe=mathgreen!50]
  Before defining any new mathematical object or concept,
  \textbf{always check whether it already exists in Mathlib}.
  Use \texttt{\#check}, \texttt{\#print}, and the
  \href{https://leanprover-community.github.io/mathlib4_docs/}{online docs}.
\end{tcolorbox}

\subsection{Key Lean Syntax Reference}

\begin{center}
\begin{tabular}{@{}lll@{}}
  \toprule
  \textbf{Symbol} & \textbf{Meaning} & \textbf{Example} \\
  \midrule
  \texttt{:}  & ``has type'' & \texttt{easy : 2 + 2 = 4} \\
  \texttt{:=} & ``is defined as'' & \texttt{:= rfl} \\
  \texttt{=}  & mathematical equality (a \texttt{Prop}) & \texttt{2 + 2 = 4} \\
  \texttt{def} & name a type or function & \texttt{def FermatLittleTheorem := \ldots} \\
  \texttt{structure} & define a multi-field object & domains, norms, etc. \\
  \texttt{class} & shared properties across types & normed spaces, etc. \\
  \bottomrule
\end{tabular}
\end{center}

\subsection{Proof Tools}

\begin{center}
\begin{tabular}{@{}lll@{}}
  \toprule
  \textbf{Tool} & \textbf{Role} & \textbf{Analogy} \\
  \midrule
  \texttt{by}   & enter tactic mode     & ``let me work step by step'' \\
  \texttt{rfl}  & prove $a = a$         & definitional equality \\
  \texttt{rw}   & rewrite with a lemma  & substitute equals for equals \\
  \texttt{calc} & chain of equalities   & showing your work line by line \\
  \texttt{sorry}& placeholder           & mark incomplete proofs \\
  \bottomrule
\end{tabular}
\end{center}

% ══════════════════════════════════════════════════════════════════════════════
\section{Mathlib Dependency Tracking}
% ══════════════════════════════════════════════════════════════════════════════

This section tracks which objects are taken from Mathlib and which are newly
defined in this project. The key principle is: \textbf{never redefine what
already exists in Mathlib}.

\subsection{Objects from Mathlib (used directly)}

\begin{center}
\begin{tabular}{@{}lll@{}}
  \toprule
  \textbf{Mathematical object} & \textbf{Mathlib name} & \textbf{Used in} \\
  \midrule
  $\mathbb{R}^n$ & \texttt{EuclideanSpace ℝ (Fin n)} & all files \\
  $C^k$ functions & \texttt{ContDiff ℝ k} & all files \\
  $C^k$ on a set & \texttt{ContDiffOn ℝ k} & \texttt{Laplace.lean} \\
  Differentiability & \texttt{Differentiable ℝ} & \texttt{Transport.lean} \\
  Differentiability at a point & \texttt{DifferentiableAt ℝ} & \texttt{Laplace.lean} \\
  Fréchet derivative & \texttt{fderiv} & \texttt{Calculus.lean} \\
  Gradient vector & \texttt{gradient} & \texttt{Calculus.lean} \\
  Scalar derivative & \texttt{deriv} & \texttt{Calculus.lean} \\
  Laplacian $\Delta$ & \texttt{Laplacian.laplacian} (\texttt{Δ}) & \texttt{Laplace.lean} \\
  Inner product $\langle\cdot,\cdot\rangle$ & \texttt{inner} / \texttt{⟪·,·⟫\_ℝ} & \texttt{Transport.lean} \\
  Inner product linear map & \texttt{innerSL ℝ} & \texttt{Laplace.lean} \\
  Continuous linear map & \texttt{ContinuousLinearMap} (\texttt{→L[ℝ]}) & \texttt{Transport.lean} \\
  $L^p$ space & \texttt{MeasureTheory.Lp} & (future) \\
  Lebesgue integral & \texttt{MeasureTheory.integral} & \texttt{Transport.lean} \\
  Interval integral & \texttt{intervalIntegral} (\texttt{∫ s in a..b}) & \texttt{Transport.lean} \\
  Average integral & \texttt{MeasureTheory.average} (\texttt{average}) & \texttt{Laplace.lean} \\
  Hausdorff measure & \texttt{Measure.hausdorffMeasure} & \texttt{Laplace.lean} \\
  Volume of unit ball & \texttt{EuclideanSpace.volume\_ball} & \texttt{Laplace.lean} \\
  Open / closed ball & \texttt{Metric.ball}, \texttt{Metric.closedBall} & \texttt{Laplace.lean} \\
  Sphere & \texttt{Metric.sphere} & \texttt{Laplace.lean} \\
  Frontier & \texttt{frontier} & \texttt{Laplace.lean} \\
  Compact support & \texttt{HasCompactSupport} & \texttt{Laplace.lean} \\
  Connected set & \texttt{IsConnected} & \texttt{Laplace.lean} \\
  Bounded set & \texttt{Bornology.IsBounded} & \texttt{Laplace.lean} \\
  \bottomrule
\end{tabular}
\end{center}

\subsection{Objects newly defined in this project}

\begin{center}
\begin{tabular}{@{}lll@{}}
  \toprule
  \textbf{Name} & \textbf{Description} & \textbf{File} \\
  \midrule
  \texttt{spatialGradient}          & $Du(x,t)$: gradient of $x'\mapsto u(x',t)$ at $x$  & \texttt{Calculus.lean} \\
  \texttt{timeDerivative}           & $u_t(x,t)$: derivative of $t'\mapsto u(x,t')$ at $t$ & \texttt{Calculus.lean} \\
  \texttt{spatialLaplacian}         & $\Delta_x u(x,t)$ for spacetime functions           & \texttt{Calculus.lean} \\
  \texttt{IsTransportSolution}      & predicate: $u_t + b \cdot Du = 0$                   & \texttt{Transport.lean} \\
  \texttt{IsInhomTransportSolution} & predicate: $u_t + b \cdot Du = f$                   & \texttt{Transport.lean} \\
  \texttt{charFlow}                 & characteristic flow $(x,t) \mapsto x - tb$          & \texttt{Transport.lean} \\
  \texttt{evansFormula}             & solution $u(x,t) = g(x-tb)$                         & \texttt{Transport.lean} \\
  \texttt{duhamelFormula}           & Duhamel solution with integral term                 & \texttt{Transport.lean} \\
  \texttt{IsHarmonic}               & predicate: $\Delta u = 0$ on $U$                    & \texttt{Laplace.lean} \\
  \texttt{IsPoissonSolution}        & predicate: $-\Delta u = f$ on $U$                   & \texttt{Laplace.lean} \\
  \texttt{fundamentalSolution}      & $\Phi(x)$, all dimensions                           & \texttt{Laplace.lean} \\
  \texttt{newtonianPotential}       & $\int \Phi(x-y) f(y)\,dy$                           & \texttt{Laplace.lean} \\
  \texttt{sphereMean}               & average of $u$ over $\partial B(x,r)$               & \texttt{Laplace.lean} \\
  \texttt{ballMean}                 & average of $u$ over $B(x,r)$                        & \texttt{Laplace.lean} \\
  \bottomrule
\end{tabular}
\end{center}

\subsection{Definitions removed after review}

The following were initially defined but later found to be redundant:

\begin{center}
\begin{tabular}{@{}lll@{}}
  \toprule
  \textbf{Name} & \textbf{Reason removed} & \textbf{Replacement} \\
  \midrule
  \texttt{laplacian} & Duplicate of \texttt{Laplacian.laplacian} & \texttt{Δ} (Mathlib) \\
  \texttt{unitBallVol} & Never used; Mathlib has \texttt{EuclideanSpace.volume\_ball} & \texttt{EuclideanSpace.volume\_ball} \\
  \texttt{laplacian\_eq\_mathlib} & Only needed due to duplicate \texttt{laplacian} & deleted \\
  \texttt{laplacian\_const\_mul} & Duplicate of \texttt{InnerProductSpace.laplacian\_smul} & Mathlib directly \\
  \texttt{fderiv\_spacetime} & Bridge lemma no longer needed & deleted \\
  \texttt{fderiv\_transport\_dir} & Bridge lemma no longer needed & deleted \\
  \texttt{inner\_spatialGradient} & Bridge lemma no longer needed & deleted \\
  \bottomrule
\end{tabular}
\end{center}

\subsection{Objects not yet in Mathlib (future work)}

\begin{center}
\begin{tabular}{@{}lll@{}}
  \toprule
  \textbf{Object} & \textbf{Needed for} & \textbf{Notes} \\
  \midrule
  Stokes theorem on spherical domains
    & Green's identity, Poisson representation
    & Main blocker; Mathlib only has rectangular domains \\
  Green's second identity on annular domains
    & Poisson representation formula
    & Depends on Stokes \\
  $n$-dimensional polar/spherical coordinates
    & $\int_{B(0,1)} \|x\|^p\,dx < \infty$, integrability of $\Phi$
    & \texttt{polarCoord} exists only for $\mathbb{R}^2$ \\
  $L^1$ integrability of $\|x\|^p$ on balls
    & \texttt{fundamentalSolution\_near\_integral\_tendsto\_zero}
    & Needs polar coordinates or dedicated Mathlib lemma \\
  \texttt{integrableOn\_rpow\_norm}
    & $\int_{B(0,r)} \|x\|^s\,dx < \infty$ for $s > -n$
    & Not in Mathlib; would unblock several proofs \\
  Sphere Hausdorff measure formula
    & $\sigma(\partial B(0,r)) = n\omega_n r^{n-1}$
    & \texttt{hausdorffMeasure\_sphere} not in Mathlib;
      blocks \texttt{fundamentalSolution\_totalFlux} \\
  Translation of set integrals
    & \texttt{fundamentalSolution\_near\_integral\_tendsto\_zero}
    & Proved manually via \texttt{integral\_sub\_left\_eq\_self} \\
  Harmonic function predicate
    & (currently defined manually as \texttt{IsHarmonic})
    & Could be added to Mathlib \\
  Weak derivative $D^\alpha u$               & Sobolev spaces & \\
  Sobolev space $W^{k,p}(\Omega)$            & Evans Ch.~5 & \\
  Sobolev norm $\|\cdot\|_{W^{k,p}}$         & Evans Ch.~5 & \\
  Trace operator $u|_{\partial\Omega}$       & boundary value problems & \\
  \bottomrule
\end{tabular}
\end{center}

% ══════════════════════════════════════════════════════════════════════════════
\section{Chapter 1: Transport Equation (Evans §2.1)}
% ══════════════════════════════════════════════════════════════════════════════

\subsection{§2.1.1 Homogeneous Case}

The homogeneous transport equation is:
\begin{equation}
  u_t + b \cdot Du = 0 \quad \text{in } \mathbb{R}^n \times (0, \infty),
  \label{eq:transport}
\end{equation}
with initial condition $u = g$ on $\mathbb{R}^n \times \{t = 0\}$.
Here $b \in \mathbb{R}^n$ is a constant vector and $g : \mathbb{R}^n \to \mathbb{R}$
is given initial data. By the method of characteristics, the unique $C^1$ solution is:
\[
  u(x, t) = g(x - tb).
\]

\subsubsection{Definitions}

The PDE predicate uses \texttt{spatialGradient} and \texttt{timeDerivative}
from \texttt{Calculus.lean}, matching Evans' notation directly:

\begin{lstlisting}[caption={PDE predicate}]
def IsTransportSolution (b : Rn) (u : Rn x R -> R) : Prop :=
  ForAll p : Rn x R,
    timeDerivative u p + <spatialGradient u p, b>_R = 0
\end{lstlisting}

The characteristic flow $(x,t) \mapsto x - tb$ is a continuous linear map.
Two simp lemmas are registered so that Lean can compute with it automatically:

\begin{lstlisting}[caption={Characteristic flow}]
noncomputable def charFlow (b : Rn) : Rn x R ->L[R] Rn :=
  ContinuousLinearMap.fst R Rn R -
  (ContinuousLinearMap.snd R Rn R).smulRight b

@[simp]
lemma charFlow_apply (b x : Rn) (t : R) :
    charFlow b (x, t) = x - t * b := by simp [charFlow]

lemma charFlow_direction_zero (b : Rn) :
    charFlow b (b, 1) = 0 := by simp
\end{lstlisting}

The solution formula is the composition of $g$ with the characteristic flow:

\begin{lstlisting}[caption={Evans formula}]
noncomputable def evansFormula (b : Rn) (g : Rn -> R) : Rn x R -> R :=
  g o charFlow b

@[simp]
lemma evansFormula_apply (b : Rn) (g : Rn -> R) (x : Rn) (t : R) :
    evansFormula b g (x, t) = g (x - t * b) := by simp [evansFormula]
\end{lstlisting}

The \texttt{@[simp]} attribute registers these as rewrite rules, so \texttt{simp}
can unfold \texttt{charFlow} and \texttt{evansFormula} automatically in proofs.

\subsubsection{Results}

\begin{lstlisting}[caption={Initial condition and regularity (proved)}]
theorem evansFormula_initial (b : Rn) (g : Rn -> R) (x : Rn) :
    evansFormula b g (x, 0) = g x

theorem evansFormula_differentiable (b : Rn) (g : Rn -> R)
    (hg : Differentiable R g) :
    Differentiable R (evansFormula b g)
\end{lstlisting}

\begin{proof}[Proof of \texttt{evansFormula\_initial}]
By definition, $u(x,0) = g(\texttt{charFlow}\,b\,(x,0)) = g(x - 0 \cdot b) = g(x)$.
\end{proof}

\begin{proof}[Proof of \texttt{evansFormula\_differentiable}]
Since $u = g \circ \texttt{charFlow}\,b$ and \texttt{charFlow} $b$ is a continuous
linear map, it is differentiable. By the chain rule, the composition of two
differentiable maps is differentiable. Hence $u$ is differentiable whenever $g$ is.
\end{proof}

\begin{lstlisting}[caption={Spatial gradient pulls back along charFlow (sorry)}]
lemma spatialGradient_evansFormula (b : Rn) (g : Rn -> R)
    (hg : Differentiable R g) (p : Rn x R) :
    spatialGradient (evansFormula b g) p = gradient g (charFlow b p) := by
  simp only [spatialGradient, evansFormula, Function.comp, charFlow_apply]
  sorry
\end{lstlisting}

\begin{lstlisting}[caption={Time derivative equals minus inner product with b (sorry)}]
lemma timeDerivative_evansFormula (b : Rn) (g : Rn -> R)
    (hg : Differentiable R g) (p : Rn x R) :
    timeDerivative (evansFormula b g) p = -<gradient g (charFlow b p), b>_R := by
  simp only [timeDerivative, evansFormula, Function.comp, charFlow_apply]
  sorry
\end{lstlisting}

Both follow from the chain rule: $x \mapsto x - tb$ has derivative $\text{id}$,
so the spatial gradient of $g(x-tb)$ is just $\nabla g(x-tb)$; and
$\partial_t[g(x-tb)] = \nabla g(x-tb) \cdot (-b) = -\langle \nabla g(x-tb), b \rangle$.
Once these are proved, the existence theorem closes by \texttt{simp}:

\begin{lstlisting}[caption={Existence (proved modulo two chain rule lemmas)}]
theorem evansFormula_solves_transport (b : Rn) (g : Rn -> R)
    (hg : Differentiable R g) :
    IsTransportSolution b (evansFormula b g) := by
  intro p
  rw [timeDerivative_evansFormula b g hg p,
      spatialGradient_evansFormula b g hg p]
  simp [real_inner_comm]
\end{lstlisting}

\begin{lstlisting}[caption={Uniqueness via characteristics (sorry)}]
theorem evansFormula_unique (b : Rn) (g : Rn -> R)
    (u : Rn x R -> R)
    (hu_pde  : IsTransportSolution b u)
    (hu_init : ForAll x : Rn, u (x, 0) = g x)
    (hu_diff : Differentiable R u) :
    u = evansFormula b g := by
  -- proved: see Transport.lean for full proof
  -- key: z(s) = u(x+sb, t+s) has z' = 0 by chain rule + PDE,
  -- so z is constant; z(-t) = g(x-tb) gives u(x,t) = g(x-tb)
\end{lstlisting}

The uniqueness proof is fully formalized. The key steps are:
\begin{enumerate}
  \item \textbf{\texttt{const\_of\_deriv\_zero}}: a function $f : \mathbb{R} \to \mathbb{R}$
    with $f' \equiv 0$ is constant, proved via
    \texttt{image\_sub\_le\_mul\_sub\_of\_deriv\_le} (MVT).
    Note: \texttt{isConst\_of\_deriv\_eq\_zero} does not exist in Mathlib.
  \item \textbf{Characteristic path}: $\gamma(s) = (x + sb, t+s)$ has derivative
    $(b, 1)$ by \texttt{HasDerivAt.prodMk}.
  \item \textbf{Chain rule}: $z'(s) = \text{fderiv}\,u\,(\gamma(s))\,(b,1)$,
    which equals $\langle \text{spatialGradient}\,u, b\rangle + \text{timeDerivative}\,u = 0$
    by the PDE hypothesis (inline of \texttt{fderiv\_transport\_dir}).
  \item \textbf{Conclusion}: $z$ is constant, $z(-t) = u(x-tb, 0) = g(x-tb)$,
    so $u(x,t) = z(0) = g(x-tb)$.
\end{enumerate}

\begin{center}
\begin{tabular}{@{}llc@{}}
  \toprule
  \textbf{Result} & \textbf{Description} & \textbf{Status} \\
  \midrule
  \texttt{evansFormula\_initial}              & $u(x,0) = g(x)$                        & \textcolor{mathgreen}{\checkmark} \\
  \texttt{evansFormula\_differentiable}       & regularity from $g$                     & \textcolor{mathgreen}{\checkmark} \\
  \texttt{spatialGradient\_evansFormula}      & chain rule for spatial gradient         & \textcolor{mathgreen}{\checkmark} \\
  \texttt{timeDerivative\_evansFormula}       & chain rule for time derivative          & \textcolor{mathgreen}{\checkmark} \\
  \texttt{evansFormula\_solves\_transport}    & existence of solution                   & \textcolor{mathgreen}{\checkmark} \\
  \texttt{evansFormula\_unique}               & uniqueness via characteristics          & \textcolor{mathgreen}{\checkmark} \\
  \bottomrule
\end{tabular}
\end{center}

\subsection{§2.1.2 Inhomogeneous Case}

The inhomogeneous transport equation introduces a source term $f$:
\begin{equation}
  u_t + b \cdot Du = f \quad \text{in } \mathbb{R}^n \times (0, \infty),
  \label{eq:inhomog_transport}
\end{equation}
with initial condition $u = g$ on $\mathbb{R}^n \times \{t = 0\}$.
By Duhamel's principle, the solution is:
\begin{equation}
  u(x, t) = g(x - tb) + \int_0^t f(x - (t-s)b,\ s)\, ds.
  \label{eq:duhamel}
\end{equation}
The first term solves the homogeneous equation; the integral corrects for the source.

\subsubsection{Definitions}

\begin{lstlisting}[caption={Inhomogeneous PDE predicate}]
def IsInhomTransportSolution (b : Rn) (f : Rn x R -> R)
    (u : Rn x R -> R) : Prop :=
  ForAll p : Rn x R,
    timeDerivative u p + <spatialGradient u p, b>_R = f p
\end{lstlisting}

\begin{lstlisting}[caption={Duhamel's formula}]
noncomputable def duhamelFormula (b : Rn) (g : Rn -> R)
    (f : Rn x R -> R) : Rn x R -> R :=
  fun p => g (p.1 - p.2 * b) +
           integral s in 0..p.2, f (p.1 - (p.2 - s) * b, s)
\end{lstlisting}

The integral uses Mathlib's interval integral \texttt{intervalIntegral}
(\texttt{∫ s in a..b, ...}).

\subsubsection{Results}

\begin{lstlisting}[caption={Initial condition: integral vanishes at t=0}]
theorem duhamelFormula_initial (b : Rn) (g : Rn -> R)
    (f : Rn x R -> R) (x : Rn) :
    duhamelFormula b g f (x, 0) = g x
\end{lstlisting}

\begin{proof}
At $t = 0$, the integral term $\int_0^0 f(x-(0-s)b, s)\,ds = 0$ since the
integration domain is empty. The remaining term is $g(x - 0 \cdot b) = g(x)$.
\end{proof}

\begin{lstlisting}[caption={Existence (partial proof)}]
theorem duhamelFormula_solves (b : Rn) (g : Rn -> R)
    (f : Rn x R -> R) (hg : Differentiable R g)
    (hf : ContDiff R 1 f) :
    IsInhomTransportSolution b f (duhamelFormula b g f) := by
  -- Step 1 (proved): partial_t[f(x-(t-s)b,s)] = -<grad f, b> via chain rule
  -- Step 2 (sorry): Leibniz rule -- Mathlib gap
  -- Step 3 (proved): time derivative of g(x-tb) via chain rule
  -- Step 4 (proved): hg_deriv + hleibniz via HasDerivAt.add
  -- Step 5 (proved): spatial gradient via
  --   hasFDerivAt_integral_of_dominated_loc_of_lip
  --   + intervalIntegral_comp_comm (sorry: measurability/Lipschitz conditions)
  -- Step 6 (proved): <Int grad f, b> = Int <grad f, b>
  --   via ContinuousLinearMap.intervalIntegral_comp_comm
  sorry -- remaining: hleibniz (Mathlib gap) + LipschitzOnWith bound condition
\end{lstlisting}

\textbf{Proof sketch:} Split $u = v + w$ where $v(x,t) = g(x-tb)$ and
$w(x,t) = \int_0^t f(x-(t-s)b, s)\,ds$. We know $v_t + b \cdot Dv = 0$.
For $w$, the Leibniz rule gives a boundary term $f(x,t)$ from the fundamental
theorem of calculus, while $\partial_t[f(x-(t-s)b,s)] = -b \cdot \nabla f$
cancels with the spatial part $b \cdot Dw$. Hence $w_t + b \cdot Dw = f$.

The proof structure is complete: all six steps are implemented, with sorry only
for technical conditions. Remaining sorry:
The proof structure is complete and most technical conditions are now discharged.
The continuity of $s \mapsto \nabla_x f(x-(t-s)b, s)$ (needed for the
\texttt{IntervalIntegrable} hypotheses and the \texttt{toDual}/\texttt{innerSL}
commutations in Steps 5--6) is proved via the identity
$\text{fderiv}\,(\lambda y.\,f(y,s))\,z = \text{fderiv}\,f\,(z,s) \circ \text{inl}$
combined with \texttt{ContDiff.continuous\_fderiv}. Two sorry remain:
\begin{itemize}
  \item \textbf{Step 2} (\texttt{hleibniz}): Mathlib has
    \texttt{integral\_hasDerivAt\_right} (FTC) and
    \texttt{hasDerivAt\_integral\_of\_dominated\_loc\_of\_lip}
    (differentiation under integral sign) separately, but no combined
    Leibniz formula for $\frac{d}{dt}\int_0^t H(t,s)\,ds = H(t,t) + \int_0^t \partial_t H(t,s)\,ds$.
    A placeholder \texttt{leibniz\_integral} is defined in \texttt{Calculus.lean}.
    This is a genuine library gap.
  \item \textbf{Step 5} (\texttt{LipschitzOnWith} hypothesis): the dominated-convergence
    bound in \texttt{hasFDerivAt\_integral\_of\_dominated\_loc\_of\_lip} requires a
    Lipschitz constant valid on the whole ball $B(x,1)$, whereas the natural bound
    $\|\text{fderiv}\,f\,(x-(t-s)b,s)\|\cdot\|b\|$ only uses the center point.
    A correct proof needs a uniform bound $\sup_{x' \in \overline{B}(x,1)}
    \|\text{fderiv}\,f\,(x'-(t-s)b,s)\|$ over the compact product
    $\overline{B}(x,1) \times [0,t]$ (via \texttt{IsCompact.exists\_bound\_of\_continuousOn}),
    which is feasible but requires substantial additional API work.
\end{itemize}
All other conditions (\texttt{AEStronglyMeasurable} of $F$ and $F'$,
\texttt{IntervalIntegrable} of the bound, the \texttt{toDual} and \texttt{innerSL}
commutations) are fully proved.

\begin{center}
\begin{tabular}{@{}llc@{}}
  \toprule
  \textbf{Result} & \textbf{Description} & \textbf{Status} \\
  \midrule
  \texttt{duhamelFormula\_initial} & $u(x,0) = g(x)$        & \textcolor{mathgreen}{\checkmark} \\
  \texttt{duhamelFormula\_solves}  & existence via Duhamel  & \textcolor{orange}{partial} \\
  \bottomrule
\end{tabular}
\end{center}

% ══════════════════════════════════════════════════════════════════════════════
\section{Chapter 2: Laplace and Poisson Equations (Evans §2.2)}
% ══════════════════════════════════════════════════════════════════════════════

\subsection{The PDEs}

The two equations covered in this section are:
\begin{align}
  -\Delta u &= 0 \quad \text{in } U \subseteq \mathbb{R}^n \quad \text{(Laplace)} \\
  -\Delta u &= f \quad \text{in } U \subseteq \mathbb{R}^n \quad \text{(Poisson)}
\end{align}
where $U$ is an open set. A solution to Laplace's equation is called \textbf{harmonic}.
The Laplacian $\Delta$ is Mathlib's \texttt{Laplacian.laplacian}, opened as \texttt{Δ}.

\begin{lstlisting}[caption={PDE predicates}]
def IsHarmonic (U : Set Rn) (u : Rn -> R) : Prop :=
  ForAll x in U, Delta u x = 0

def IsPoissonSolution (U : Set Rn) (f : Rn -> R) (u : Rn -> R) : Prop :=
  ForAll x in U, -Delta u x = f x
\end{lstlisting}

\begin{lstlisting}[caption={Laplace is Poisson with f=0 (proved)}]
lemma isHarmonic_iff_isPoissonSolution_zero (U : Set Rn) (u : Rn -> R) :
    IsHarmonic U u <-> IsPoissonSolution U 0 u
\end{lstlisting}

\begin{proof}
$u$ is harmonic if and only if $\Delta u = 0$ on $U$, which holds if and only if
$-\Delta u = 0 = f$ on $U$, i.e., $u$ solves Poisson's equation with $f = 0$.
\end{proof}

\subsection{§2.2.1 Fundamental Solution}

\begin{definition}
The \textbf{fundamental solution} $\Phi : \mathbb{R}^n \setminus \{0\} \to \mathbb{R}$
satisfies $-\Delta \Phi = \delta_0$ in the distributional sense, where $\omega_n$
denotes the volume of the unit ball in $\mathbb{R}^n$:
\[
  \Phi(x) = \begin{cases}
    \dfrac{1}{n(n-2)\omega_n} |x|^{2-n} & n \geq 3 \\[6pt]
    -\dfrac{1}{2\pi} \log|x| & n = 2
  \end{cases}
\]
\end{definition}

\begin{lstlisting}[caption={Fundamental solution smooth off zero (proved)}]
lemma fundamentalSolution_contDiff_off_zero :
    ContDiffOn R Top fundamentalSolution ({0} : Set Rn)^c
\end{lstlisting}

\begin{proof}
We check each case separately. For $n = 0$, $\Phi = 0$ is trivially smooth.
For $n = 1$, $\Phi(x) = \frac{1}{2}\|x\|$, which is smooth since $\|x\|$ is
$C^\infty$ on $\mathbb{R} \setminus \{0\}$.
For $n = 2$, $\Phi(x) = c\log\|x\|$. Since $\|x\| > 0$ for $x \neq 0$, and
$\log$ is $C^\infty$ on $(0,\infty)$, the composition is $C^\infty$.
For $n \geq 3$, $\Phi(x) = c\|x\|^{2-n}$. Since $\|x\| \neq 0$ for $x \neq 0$,
and $t \mapsto t^{2-n}$ is $C^\infty$ on $(0,\infty)$, the composition is $C^\infty$.
\end{proof}

The \textbf{Newtonian potential} is the convolution of $\Phi$ with a source $f$:

\begin{lstlisting}[caption={Newtonian potential}]
noncomputable def newtonianPotential (f : Rn -> R) : Rn -> R :=
  fun x => integral y, fundamentalSolution (x - y) * f y
\end{lstlisting}

\subsection{§2.2.1 Harmonicity of the Fundamental Solution}

This subsection contains the key analytic computation: proving $\Delta\Phi = 0$
away from the origin. The proof requires three helper lemmas.

\subsubsection{Helper definitions}

To avoid type ambiguity with Mathlib's conjugate-linear \texttt{innerSL}, we
introduce wrapper definitions with unambiguous types:

\begin{lstlisting}[caption={Real inner product CLMs (avoiding conjugate-linear ambiguity)}]
noncomputable def realInnerBiL : Rn ->L[R] Rn ->L[R] R :=
  (innerSL R : Rn ->L[R] Rn ->L[R] R)

noncomputable def realInnerL (x : Rn) : Rn ->L[R] R :=
  realInnerBiL x

lemma realInnerL_apply (x y : Rn) : realInnerL x y = <x, y>_R
\end{lstlisting}

\subsubsection{Helper lemmas (all proved)}

\begin{lstlisting}[caption={Scalar linearity of Laplacian (proved)}]
lemma laplacian_const_mul (c : R) (f : Rn -> R) (hf : ContDiffAt R 2 f x) :
    Delta (fun y => c * f y) x = c * Delta f x
\end{lstlisting}

\begin{proof}
Since $\Delta$ is a linear operator, $\Delta(c \cdot f) = c \cdot \Delta f$
for any constant $c \in \mathbb{R}$ and $C^2$ function $f$.
\end{proof}

\begin{lstlisting}[caption={First derivative of radial power (proved)}]
lemma hasFDerivAt_norm_rpow_of_ne (x : Rn) (hx : x != 0) (p : R) :
    HasFDerivAt (fun x => norm x ^ p)
      ((p * norm x ^ (p - 2)) * realInnerL x) x
\end{lstlisting}

\begin{proof}
Write $\|x\|^p = (\|x\|^2)^{p/2}$. By the chain rule:
\[
  D(\|x\|^p)(x)(v) = \frac{p}{2}(\|x\|^2)^{p/2 - 1} \cdot D(\|x\|^2)(x)(v)
  = \frac{p}{2}\|x\|^{p-2} \cdot 2\langle x, v\rangle = p\|x\|^{p-2}\langle x, v\rangle.
\]
Hence the Fréchet derivative at $x$ is the linear map $v \mapsto p\|x\|^{p-2}\langle x, v\rangle$.
\end{proof}

\begin{lstlisting}[caption={Laplacian of radial power (proved)}]
lemma laplacian_norm_rpow_eq (p : R) (x : Rn) (hx : x != 0) :
    Delta (fun x => norm x ^ p) x = p * (n + p - 2) * norm x ^ (p - 2)
\end{lstlisting}

\begin{proof}
From \texttt{hasFDerivAt\_norm\_rpow\_of\_ne}, the gradient is
$\nabla(\|x\|^p) = p\|x\|^{p-2}x$.
Differentiating again using the product rule:
\[
  D^2(\|x\|^p)(x)(v,w) = p(p-2)\|x\|^{p-4}\langle x,v\rangle\langle x,w\rangle
  + p\|x\|^{p-2}\langle v,w\rangle.
\]
Taking the trace over an orthonormal basis $\{e_i\}$:
\begin{align*}
  \Delta(\|x\|^p)(x)
  &= \sum_i D^2(\|x\|^p)(x)(e_i, e_i) \\
  &= p(p-2)\|x\|^{p-4}\underbrace{\sum_i\langle x,e_i\rangle^2}_{=\|x\|^2}
   + p\|x\|^{p-2}\underbrace{\sum_i\langle e_i,e_i\rangle}_{=n} \\
  &= p(p-2)\|x\|^{p-2} + pn\|x\|^{p-2} = p(n+p-2)\|x\|^{p-2}.
\end{align*}
\end{proof}

\begin{lstlisting}[caption={Laplacian of log norm (proved)}]
lemma laplacian_log_norm_eq (x : Rn) (hx : x != 0) :
    Delta (fun x => log (norm x)) x = (n - 2) * norm x ^ (-2)
\end{lstlisting}

\begin{proof}
The proof follows the same structure as \texttt{laplacian\_norm\_rpow\_eq}.

\textbf{Step 1: First derivative.}
Write $\log\|y\| = \log(\|y\|^1)$ and apply the chain rule.
Since $D(\|y\|)(y)(v) = \|y\|^{-1}\langle y, v\rangle$, and
$\frac{d}{dt}\log t\big|_{t=\|y\|} = \|y\|^{-1}$, the chain rule gives:
\[
  D(\log\|y\|)(y)(v) = \|y\|^{-1} \cdot \|y\|^{-1}\langle y, v\rangle
  = \|y\|^{-2}\langle y, v\rangle.
\]
Hence $\nabla(\log\|y\|) = \|y\|^{-2}y$, and the first derivative is
the linear map $v \mapsto \|y\|^{-2}\langle y, v\rangle$.

\textbf{Step 2: Second derivative.}
Differentiating $y \mapsto \|y\|^{-2}\langle y, v\rangle$ using the product rule:
\[
  D^2(\log\|x\|)(x)(v,w)
  = -2\|x\|^{-4}\langle x,v\rangle\langle x,w\rangle + \|x\|^{-2}\langle v,w\rangle.
\]

\textbf{Step 3: Trace over orthonormal basis.}
Taking the trace over an orthonormal basis $\{e_i\}$:
\begin{align*}
  \Delta(\log\|x\|)
  &= \sum_i D^2(\log\|x\|)(x)(e_i,e_i) \\
  &= -2\|x\|^{-4}\underbrace{\sum_i\langle x,e_i\rangle^2}_{=\|x\|^2}
   + \|x\|^{-2}\underbrace{\sum_i\langle e_i,e_i\rangle}_{=n} \\
  &= -2\|x\|^{-2} + n\|x\|^{-2} = (n-2)\|x\|^{-2}.
\end{align*}
In particular, when $n = 2$, $\Delta(\log\|x\|) = 0$, confirming that
$\log\|x\|$ is harmonic in $\mathbb{R}^2 \setminus \{0\}$.
\end{proof}

\subsubsection{Main result}

\begin{lstlisting}[caption={Fundamental solution is harmonic off zero (proved modulo log lemma)}]
lemma fundamentalSolution_harmonic_off_zero (x : Rn) (hx : x != 0) :
    Delta fundamentalSolution x = 0
\end{lstlisting}

\begin{proof}
We consider each dimension separately.

\textbf{Case $n = 0$:} $\Phi \equiv 0$, so $\Delta\Phi = 0$ trivially.

\textbf{Case $n = 1$:} $\Phi(x) = \frac{1}{2}\|x\|^1$.
By \texttt{laplacian\_const\_mul} and \texttt{laplacian\_norm\_rpow\_eq} with $p = 1$:
\[ \Delta\Phi = \frac{1}{2} \cdot 1 \cdot (1 + 1 - 2) \cdot \|x\|^{-1} = 0. \]

\textbf{Case $n = 2$:} $\Phi(x) = -\frac{1}{2\pi}\log\|x\|$.
By \texttt{laplacian\_const\_mul} and \texttt{laplacian\_log\_norm\_eq}:
\[ \Delta\Phi = -\frac{1}{2\pi} \cdot (2-2) \cdot \|x\|^{-2} = 0. \]

\textbf{Case $n \geq 3$:} $\Phi(x) = \frac{1}{n(n-2)\omega_n}\|x\|^{2-n}$.
By \texttt{laplacian\_const\_mul} and \texttt{laplacian\_norm\_rpow\_eq} with $p = 2-n$:
\[ \Delta\Phi = c \cdot (2-n) \cdot (n + (2-n) - 2) \cdot \|x\|^{-n}
             = c \cdot (2-n) \cdot 0 = 0. \]

In all cases $\Delta\Phi(x) = 0$ for $x \neq 0$.
\end{proof}

\subsection{§2.2.2 Mean Value Property}

\begin{lstlisting}[caption={Mean value property (sorry)}]
theorem harmonic_sphereMeanValue ... : u x = sphereMean u x r := by sorry
theorem harmonic_ballMeanValue   ... : u x = ballMean u x r  := by sorry
theorem meanValue_implies_harmonic ... : IsHarmonic U u       := by sorry
\end{lstlisting}

Proof sketch: apply Green's identity to $u$ and $\Phi(\cdot - x)$ on
$B(x,r) \setminus B(x,\varepsilon)$, then let $\varepsilon \to 0$.

\subsection{§2.2.3 Maximum Principles and Regularity}

\begin{lstlisting}[caption={Maximum principles and regularity (all sorry)}]
theorem harmonic_weakMax   ... : ForAll x in U, u x <= sSup (u '' frontier U)
theorem harmonic_strongMax ... : ForAll x in U, u x = u x0
theorem harmonic_smooth    ... : ContDiffOn R Top u U
\end{lstlisting}

\subsection{§2.2.4 Representation Formula}

The proof of the representation formula requires the following intermediate lemmas,
whose status reflects the current state of Mathlib:

\begin{center}
\begin{tabular}{@{}llc@{}}
  \toprule
  \textbf{Lemma} & \textbf{Blocker} & \textbf{Status} \\
  \midrule
  \texttt{fundamentalSolution\_near\_integral\_tendsto\_zero}
    & $\|\Phi\|$ integrable on $B(0,1)$: needs $n$-dim polar coordinates & \textcolor{orange}{partial} \\
  \texttt{green\_identity\_annulus}
    & Gauss--Green on spherical annulus: needs Stokes & \textcolor{orange}{partial} \\
  \texttt{fundamentalSolution\_totalFlux}
    & sphere Hausdorff measure $\sigma(\partial B(0,\varepsilon)) = n\omega_n\varepsilon^{n-1}$:
      \texttt{hausdorffMeasure\_sphere} not in Mathlib & \textcolor{sorryred}{sorry} \\
  \texttt{green\_boundary\_tendsto\_f}
    & depends on \texttt{green\_identity\_annulus} and \texttt{fundamentalSolution\_totalFlux} & \textcolor{sorryred}{sorry} \\
  \texttt{newtonianPotential\_solves\_poisson}
    & depends on all above & \textcolor{sorryred}{sorry} \\
  \bottomrule
\end{tabular}
\end{center}

The entire proof chain for the Poisson representation formula is blocked by three
independent missing tools in Mathlib, each of which would require substantial
infrastructure to develop:

\begin{enumerate}
  \item \textbf{Stokes' theorem on spherical domains.}
    Required by \texttt{green\_identity\_annulus} (Step 2) and
    \texttt{green\_boundary\_tendsto\_f}.
    Mathlib's divergence theorem (\texttt{integral\_divergence\_of\_hasFDerivWithinAt\_off\_countable})
    covers only rectangular boxes; the spherical case is not available.

  \item \textbf{Sphere Hausdorff measure formula.}
    Required by \texttt{fundamentalSolution\_totalFlux}:
    the computation $\int_{\partial B(0,\varepsilon)} d\sigma = n\omega_n\varepsilon^{n-1}$
    requires \texttt{hausdorffMeasure\_sphere} or equivalent, which is not in Mathlib.
    The Evans computation (§2.2.4, p.~23) gives
    $\int_{\partial B(0,\varepsilon)} \langle\nabla\Phi, \nu\rangle\,d\sigma
    = c_n(2-n)\varepsilon^{1-n} \cdot n\omega_n\varepsilon^{n-1} = -1$,
    but the surface area step is unavailable.

  \item \textbf{$n$-dimensional polar coordinates.}
    Required by \texttt{fundamentalSolution\_near\_integral\_tendsto\_zero}
    (integrability of $\|\Phi\|$ on $B(0,1)$) and implicitly by
    \texttt{fundamentalSolution\_totalFlux}.
    Mathlib provides \texttt{polarCoord} for $\mathbb{R}^2$ only;
    the general $\mathbb{R}^n$ version is not available.
\end{enumerate}

Despite these blockers, significant partial progress has been made:
\texttt{green\_identity\_annulus} Step~1 (the algebraic identity
$v\,\Delta u - u\,\Delta v = \sum_i \partial_{e_i}(v\,\partial_{e_i}u - u\,\partial_{e_i}v)$)
is fully proved, and the framework of
\texttt{fundamentalSolution\_near\_integral\_tendsto\_zero}
(translation, sequential convergence, \texttt{nhdsWithin} transfer) is complete.

\subsection{Status Summary}

\begin{center}
\begin{tabular}{@{}llc@{}}
  \toprule
  \textbf{Result} & \textbf{Description} & \textbf{Status} \\
  \midrule
  \texttt{isHarmonic\_iff\_isPoissonSolution\_zero} & Laplace is Poisson with $f=0$ & \textcolor{mathgreen}{\checkmark} \\
  \texttt{fundamentalSolution\_contDiff\_off\_zero} & $\Phi \in C^\infty(\mathbb{R}^n \setminus \{0\})$ & \textcolor{mathgreen}{\checkmark} \\
  \texttt{laplacian\_const\_mul}     & linearity of $\Delta$              & \textcolor{mathgreen}{\checkmark} \\
  \texttt{hasFDerivAt\_norm\_rpow\_of\_ne} & $D(\|x\|^p)$ formula        & \textcolor{mathgreen}{\checkmark} \\
  \texttt{laplacian\_norm\_rpow\_eq} & $\Delta(\|x\|^p) = p(n+p-2)\|x\|^{p-2}$ & \textcolor{mathgreen}{\checkmark} \\
  \texttt{laplacian\_log\_norm\_eq}  & $\Delta(\log\|x\|) = (n-2)\|x\|^{-2}$ & \textcolor{mathgreen}{\checkmark} \\
  \texttt{fundamentalSolution\_harmonic\_off\_zero} & $\Delta\Phi = 0$ off origin & \textcolor{mathgreen}{\checkmark} \\
  \texttt{integral\_ball\_translate} & $\int_{B(x,\varepsilon)} f(x-y)\,dy = \int_{B(0,\varepsilon)} f(z)\,dz$ & \textcolor{mathgreen}{\checkmark} \\
  \texttt{green\_identity\_annulus} & Green's 2nd identity on annulus (Step 1 proved) & \textcolor{orange}{partial} \\
  \texttt{fundamentalSolution\_near\_integral\_tendsto\_zero} & $\int_{B(x,\varepsilon)}\|\Phi\| \to 0$ & \textcolor{orange}{partial} \\
  \texttt{harmonic\_sphereMeanValue} & mean value on spheres              & \textcolor{sorryred}{sorry} \\
  \texttt{harmonic\_ballMeanValue}   & mean value on balls                & \textcolor{sorryred}{sorry} \\
  \texttt{meanValue\_implies\_harmonic} & converse of MVP                & \textcolor{sorryred}{sorry} \\
  \texttt{harmonic\_weakMax}         & weak maximum principle             & \textcolor{sorryred}{sorry} \\
  \texttt{harmonic\_strongMax}       & strong maximum principle           & \textcolor{sorryred}{sorry} \\
  \texttt{harmonic\_smooth}          & harmonic $\Rightarrow C^\infty$    & \textcolor{sorryred}{sorry} \\
  \texttt{newtonianPotential\_solves\_poisson} & representation formula   & \textcolor{sorryred}{sorry} \\
  \bottomrule
\end{tabular}
\end{center}

% ══════════════════════════════════════════════════════════════════════════════
\section{Chapter 3: Heat Equation (Evans §2.3)}
% ══════════════════════════════════════════════════════════════════════════════

\subsection{The PDEs}

The initial value problem for the heat (diffusion) equation:
\begin{align}
  u_t - \Delta u &= 0 \quad \text{in } \mathbb{R}^n \times (0,\infty) \quad \text{(homogeneous)} \\
  u_t - \Delta u &= f \quad \text{in } \mathbb{R}^n \times (0,\infty) \quad \text{(inhomogeneous)} \\
  u &= g \quad \text{on } \mathbb{R}^n \times \{t = 0\}.
\end{align}
Here $u_t$ is \texttt{timeDerivative} and $\Delta u$ is the spatial Laplacian
\texttt{spatialLaplacian}, both from \texttt{Calculus.lean}.

Unlike the Laplace chapter, the heat chapter avoids the three Mathlib blockers
(Stokes on spheres, \texttt{hausdorffMeasure\_sphere}, $n$-dim polar coordinates):
the fundamental solution is an \emph{explicit} Gaussian, and its normalization uses
Mathlib's multivariate Gaussian integral directly.

\subsection{§2.3.1 The Heat Kernel}

The fundamental solution (heat kernel) is
\begin{equation}
  \Phi(x, t) =
  \begin{cases}
    (4\pi t)^{-n/2} \exp\!\left(-\dfrac{|x|^2}{4t}\right) & t > 0, \\[1ex]
    0 & t \le 0.
  \end{cases}
\end{equation}

\begin{lstlisting}[caption={Heat kernel definition}]
noncomputable def heatKernel (p : Rn x R) : R :=
  if 0 < p.2 then
    (4 * pi * p.2) ^ (-n / 2) * exp (-||p.1||^2 / (4 * p.2))
  else 0

noncomputable def heatSolution (g : Rn -> R) : Rn x R -> R :=
  fun p => Int y, heatKernel (p.1 - y, p.2) * g y
\end{lstlisting}

The key computations are the time derivative and spatial Laplacian of the kernel,
which turn out to be \emph{equal}, so the kernel solves the heat equation:
\begin{align}
  \Phi_t(x,t) &= \Phi(x,t) \cdot \left(\frac{|x|^2}{4t^2} - \frac{n}{2t}\right), \\
  \Delta \Phi(x,t) &= \Phi(x,t) \cdot \left(\frac{|x|^2}{4t^2} - \frac{n}{2t}\right).
\end{align}

\textbf{Normalization} ($\int_{\mathbb{R}^n} \Phi(\cdot, t) = 1$): the kernel is
$(4\pi t)^{-n/2}$ times the Gaussian $\exp(-|x|^2/4t)$, whose integral over
$\mathbb{R}^n$ is $(4\pi t)^{n/2}$ by Mathlib's
\texttt{GaussianFourier.integral\_rexp\_neg\_mul\_sq\_norm}. The two powers cancel.

\textbf{Time derivative}: on $t > 0$ the kernel agrees with the smooth branch
$F(s) = (4\pi s)^{-n/2} \exp(-|x|^2/4s)$. Writing $\Phi = u \cdot v$ with
$u(s) = (4\pi s)^{-n/2}$, $v(s) = \exp(-|x|^2/4s)$, the product and chain rules give the
claim. The derivative is transferred to the kernel via \texttt{Filter.EventuallyEq}.

\textbf{Spatial Laplacian}: the Gaussian Laplacian is computed by expanding over the
standard orthonormal basis (\texttt{laplacian\_eq\_iteratedFDeriv\_orthonormalBasis}),
computing the second directional derivative along each $e_i$ via the product rule, and
summing using Parseval ($\sum_i \langle x, e_i \rangle^2 = |x|^2$). The constant
$(4\pi t)^{-n/2}$ pulls out of the Laplacian via \texttt{laplacian\_smul}.

\begin{center}
\begin{tabular}{@{}llc@{}}
  \toprule
  \textbf{Result} & \textbf{Description} & \textbf{Status} \\
  \midrule
  \texttt{isHeatSolution\_iff\_isInhomHeatSolution\_zero} & homogeneous $=$ source $0$ & \textcolor{mathgreen}{\checkmark} \\
  \texttt{heatKernel\_of\_nonpos} & $\Phi = 0$ for $t \le 0$ & \textcolor{mathgreen}{\checkmark} \\
  \texttt{heatKernel\_pos} & $\Phi > 0$ for $t > 0$ & \textcolor{mathgreen}{\checkmark} \\
  \texttt{heatKernel\_nonneg} & $\Phi \ge 0$ & \textcolor{mathgreen}{\checkmark} \\
  \texttt{heatKernel\_integral\_eq\_one} & $\int_{\mathbb{R}^n} \Phi(\cdot,t) = 1$ & \textcolor{mathgreen}{\checkmark} \\
  \texttt{heatKernel\_timeDerivative} & $\Phi_t = \Phi(|x|^2/4t^2 - n/2t)$ & \textcolor{mathgreen}{\checkmark} \\
  \texttt{gaussian\_laplacian} & $\Delta_y \exp(-|y|^2/4t)$ & \textcolor{mathgreen}{\checkmark} \\
  \texttt{heatKernel\_spatialLaplacian} & $\Delta \Phi = \Phi(|x|^2/4t^2 - n/2t)$ & \textcolor{mathgreen}{\checkmark} \\
  \texttt{heatKernel\_solves\_heat} & $\Phi_t - \Delta\Phi = 0$ ($t>0$) & \textcolor{mathgreen}{\checkmark} \\
  \texttt{heatSolution\_solves\_heat} & convolution solves heat eqn & \textcolor{sorryred}{sorry} \\
  \bottomrule
\end{tabular}
\end{center}

The only remaining sorry is \texttt{heatSolution\_solves\_heat}: that the convolution
$u(x,t) = \int \Phi(x-y,t)\,g(y)\,dy$ solves the heat equation. This requires
differentiation under the integral sign (in both $x$ and $t$), pushing the derivatives
onto the kernel and applying \texttt{heatKernel\_solves\_heat}, with dominated-convergence
bounds for the Gaussian and its derivatives. The kernel computations it depends on are
all complete.

% ══════════════════════════════════════════════════════════════════════════════
\section{Calculus Utilities: \texttt{Calculus.lean}}
% ══════════════════════════════════════════════════════════════════════════════

\subsection{Motivation}

\texttt{Calculus.lean} provides a shared library of spacetime calculus definitions
for functions $u : \mathbb{R}^n \times \mathbb{R} \to \mathbb{R}$, matching Evans'
notation throughout the project.

\subsection{Definitions}

\subsubsection{Spatial Gradient}

\begin{definition}
The \textbf{spatial gradient} $Du(x,t)$ is the gradient of the map
$x' \mapsto u(x', t)$ at $x$, viewed as a vector in $\mathbb{R}^n$.
\end{definition}

\begin{lstlisting}[caption={Spatial gradient in Lean}]
noncomputable def spatialGradient (u : Rn x R -> R) (p : Rn x R) : Rn :=
  gradient (fun x => u (x, p.2)) p.1
\end{lstlisting}

In Mathlib, \texttt{gradient f x} is defined for $f : F \to \mathbb{R}$ on an
inner product space $F$, and returns the unique vector $\nabla f(x) \in F$
satisfying $\langle \nabla f(x), v \rangle = \text{fderiv}_{\mathbb{R}}\, f\, x\, v$
for all $v$. Note: \texttt{gradient} does \emph{not} take $\mathbb{R}$ as an
explicit argument --- the field is inferred automatically from the function's type.

\subsubsection{Time Derivative}

\begin{definition}
The \textbf{time derivative} $u_t(x,t)$ is the derivative of the map
$t' \mapsto u(x, t')$ at $t$.
\end{definition}

\begin{lstlisting}[caption={Time derivative in Lean}]
noncomputable def timeDerivative (u : Rn x R -> R) (p : Rn x R) : R :=
  deriv (fun t => u (p.1, t)) p.2
\end{lstlisting}

Here \texttt{deriv} is Mathlib's scalar derivative for functions $\mathbb{R} \to \mathbb{R}$.
When $u$ is not differentiable, \texttt{deriv} returns $0$ by convention,
so no differentiability hypothesis is needed in the definition itself.

\subsubsection{Laplacian}

The Laplacian $\Delta u$ is \textbf{not} redefined in \texttt{Calculus.lean}.
Instead, \texttt{Laplace.lean} uses Mathlib's \texttt{Laplacian.laplacian} directly
(opened as \texttt{Δ}). The \texttt{spatialLaplacian} wrapper is provided for
spacetime functions:

\begin{lstlisting}[caption={Spatial Laplacian for spacetime functions}]
noncomputable def spatialLaplacian (u : Rn x R -> R) (p : Rn x R) : R :=
  Laplacian.laplacian (fun x => u (x, p.2)) p.1
\end{lstlisting}

\subsection{Status Summary}

\begin{center}
\begin{tabular}{@{}llc@{}}
  \toprule
  \textbf{Definition} & \textbf{Description} & \textbf{Status} \\
  \midrule
  \texttt{spatialGradient}  & $Du(x,t)$ as a vector in $\mathbb{R}^n$          & \textcolor{mathgreen}{\checkmark} \\
  \texttt{timeDerivative}   & $u_t(x,t)$ as a scalar                            & \textcolor{mathgreen}{\checkmark} \\
  \texttt{spatialLaplacian} & $\Delta_x u(x,t)$ wrapping \texttt{Δ}             & \textcolor{mathgreen}{\checkmark} \\
  \bottomrule
\end{tabular}
\end{center}

% ══════════════════════════════════════════════════════════════════════════════
\section{Upcoming: Sobolev Spaces (Evans Ch.~5)}
% ══════════════════════════════════════════════════════════════════════════════

This is the \textbf{main target} of the project. All of modern PDE theory is
phrased in the language of Sobolev spaces, and Mathlib has almost nothing here.

\subsection{Key definitions to formalize}

\begin{enumerate}
  \item \textbf{Weak derivative} (Evans Definition 5.1):
    $v \in L^1_{\mathrm{loc}}(\Omega)$ is the $\alpha$-th weak derivative of
    $u \in L^1_{\mathrm{loc}}(\Omega)$ if
    \[
      \int_\Omega u\, D^\alpha \phi \, dx
      = (-1)^{|\alpha|} \int_\Omega v\, \phi \, dx
      \quad \forall \phi \in C^\infty_c(\Omega).
    \]

  \item \textbf{Sobolev space} $W^{k,p}(\Omega)$:
    \[
      W^{k,p}(\Omega) = \{ u \in L^p(\Omega) :
      D^\alpha u \in L^p(\Omega) \text{ for all } |\alpha| \leq k \}
    \]
    with norm $\|u\|_{W^{k,p}} = \left(\sum_{|\alpha|\leq k}
    \|D^\alpha u\|_{L^p}^p\right)^{1/p}$.

  \item \textbf{Sobolev embedding theorem.}
  \item \textbf{Trace theorem.}
\end{enumerate}

\begin{todo}
  Check Mathlib for any existing infrastructure around
  \texttt{MeasureTheory.Lp} that could serve as a foundation
  for $W^{k,p}$.
\end{todo}

% ══════════════════════════════════════════════════════════════════════════════
\section{Resources}
% ══════════════════════════════════════════════════════════════════════════════

\begin{itemize}
  \item \href{https://leanprover-community.github.io/mathematics_in_lean/}
    {Mathematics in Lean} — primary tutorial
  \item \href{https://leanprover-community.github.io/mathlib4_docs/}
    {Mathlib4 Documentation} — search before defining
  \item \href{https://adam.math.hhu.de}{Natural Number Game} — interactive basics
  \item L.C. Evans, \textit{Partial Differential Equations}, AMS Graduate Studies
    in Mathematics, Vol.~19, 1998
\end{itemize}

\end{document}